% © 2025 Ing. David Jaroš — CC BY-NC-ND 4.0
%
% This work is licensed under a Creative Commons Attribution-NonCommercial-NoDerivatives
% 4.0 International License (CC BY-NC-ND 4.0).
%
% License History: Earlier drafts (up to v0.3) were released under CC BY 4.0.
% From v0.4 onward, all material is released under CC BY-NC-ND 4.0 to protect
% the integrity of the theoretical work during ongoing academic development.
%
% See LICENSE.md for full license text.

\ifdefined\INCLUDEMODE
  % Being included in another document - skip preamble
\else
  \documentclass[12pt]{article}
  \usepackage[a4paper, margin=2.5cm]{geometry}
  \usepackage{amsmath, amssymb, amsthm}
  \usepackage{hyperref}
  \usepackage{tcolorbox}

  \newtheorem{theorem}{Theorem}[section]
  \newtheorem{lemma}[theorem]{Lemma}
  \newtheorem{corollary}[theorem]{Corollary}
  \newtheorem{proposition}[theorem]{Proposition}
  \theoremstyle{definition}
  \newtheorem{definition}[theorem]{Definition}
  \theoremstyle{remark}
  \newtheorem{remark}[theorem]{Remark}

  \title{\textbf{GR Closure Step 3: Lorentzian Signature Theorem\\
         Deriving $(-,+,+,+)$ from the Algebra $\mathbb{C}\otimes\mathbb{H}$}}
  \author{Ing. David Jaroš}
  \date{March 2026}

  \begin{document}
  \maketitle

  \begin{abstract}
  The GR Recovery Audit identified that the Lorentzian signature $(-,+,+,+)$ of the
  UBT emergent metric is currently \emph{enforced by normalization choice}, not derived
  from the algebraic structure of biquaternions.  We provide a sharper statement: the
  algebra $\mathbb{C}\otimes\mathbb{H}$ contains a natural $4$-dimensional real subspace
  spanned by a canonical basis that carries a quadratic form of signature $(1,3)$
  (equivalently $(-,+,+,+)$ in the physics convention).  We then show that the
  normalization $\mathcal{N}$ merely \emph{selects} this canonical signature rather
  than imposing it arbitrarily, and we identify the residual freedom.
  \end{abstract}
\fi

\section{Lorentzian Signature from $\mathbb{C}\otimes\mathbb{H}$}
\label{sec:signature_theorem}

\begin{tcolorbox}[colback=blue!5!white,colframe=blue!75!black,title=Purpose of This Step]
\textbf{Sharpening the signature argument identified in the GR Recovery Audit.}

The audit's critique: \emph{``Signature is enforced by normalization choice,
not derived from algebraic hypotheses.''}

This step provides an algebraic derivation showing that $\mathbb{C}\otimes\mathbb{H}$
naturally carries a quadratic form of signature $(1,3)$, so the normalization $\mathcal{N}$
makes a \emph{canonical} rather than \emph{arbitrary} choice.
\end{tcolorbox}

% -----------------------------------------------------------------------
\subsection{Algebraic Structure of $\mathbb{C}\otimes\mathbb{H}$}

The biquaternion algebra is
\begin{equation}
  \mathbb{B} \;:=\; \mathbb{C}\otimes_{\mathbb{R}}\mathbb{H}
  \;\cong\; M_2(\mathbb{C})
  \quad\text{(as a $\mathbb{C}$-algebra)},
  \label{eq:biquaternion_algebra}
\end{equation}
where the isomorphism to $2\times 2$ complex matrices is given by the standard
Pauli representation:
\begin{equation}
  \mathbf{1}\mapsto I_2,\quad
  \mathbf{i}\mapsto i\sigma_3,\quad
  \mathbf{j}\mapsto i\sigma_2,\quad
  \mathbf{k}\mapsto i\sigma_1,
  \label{eq:pauli_rep}
\end{equation}
with the Pauli matrices $\sigma_1,\sigma_2,\sigma_3$.

\begin{definition}[Real-vector subspace $V_{\mathbb{R}}$]
\label{def:real_vector_subspace}
Inside $\mathbb{B}$, define the four-dimensional \emph{real-vector subspace}
\begin{equation}
  V_{\mathbb{R}} \;:=\; \mathrm{span}_{\mathbb{R}}\{1,\,i\mathbf{i},\,i\mathbf{j},\,i\mathbf{k}\}
  \;\subset\; \mathbb{B}.
  \label{eq:Vreal}
\end{equation}
Elements of $V_{\mathbb{R}}$ have the form $v = t\cdot\mathbf{1} + x\cdot i\mathbf{i}
+ y\cdot i\mathbf{j} + z\cdot i\mathbf{k}$ with $t,x,y,z\in\mathbb{R}$.
\end{definition}

\begin{remark}
In the Pauli representation, elements of $V_{\mathbb{R}}$ correspond to $2\times 2$
Hermitian matrices $v = t\,I_2 - x\,\sigma_3 - y\,\sigma_2 - z\,\sigma_1$,
i.e.\ to real linear combinations of $\{I_2,\sigma_1,\sigma_2,\sigma_3\}$
(the Hermitian basis of $M_2(\mathbb{C})$).
\end{remark}

% -----------------------------------------------------------------------
\subsection{The Natural Quadratic Form}

\begin{definition}[Biquaternionic scalar quadratic form]
\label{def:quad_form}
For $v\in V_{\mathbb{R}}$, define
\begin{equation}
  Q(v) \;:=\; \mathrm{Sc}(v\cdot v^\dagger)
  \;=\; \tfrac{1}{4}\mathrm{Tr}(v\cdot v^\dagger).
  \label{eq:quad_form}
\end{equation}
\end{definition}

\begin{lemma}[Signature of $Q$]
\label{lem:signature}
For $v = t\,\mathbf{1} + x\,(i\mathbf{i}) + y\,(i\mathbf{j}) + z\,(i\mathbf{k})
\in V_{\mathbb{R}}$, the quadratic form \eqref{eq:quad_form} evaluates to
\begin{equation}
  Q(v) \;=\; t^2 - x^2 - y^2 - z^2.
  \label{eq:minkowski_form}
\end{equation}
The associated bilinear form has signature $(1,3)$, i.e.\ $(-,+,+,+)$ in the physics
convention $\eta = \mathrm{diag}(-1,+1,+1,+1)$ upon setting $v^0=t$, $v^i = (x,y,z)$.
\end{lemma}

\begin{proof}
Compute directly using the quaternion multiplication rules
$\mathbf{i}^2=\mathbf{j}^2=\mathbf{k}^2=-1$ and
$\mathbf{i}\mathbf{j}=\mathbf{k}$, $\mathbf{j}\mathbf{k}=\mathbf{i}$, $\mathbf{k}\mathbf{i}=\mathbf{j}$.

For the adjoint, $\mathbf{1}^\dagger=\mathbf{1}$ and $(i\mathbf{e}_a)^\dagger = -i\mathbf{e}_a$
(complex conjugation reverses $i$, and quaternionic conjugation reverses $\mathbf{e}_a$,
giving two minus signs that cancel; however the combined biquaternionic adjoint gives
$(i\mathbf{e}_a)^\dagger = \overline{i}\,(-\mathbf{e}_a) = (-i)(-\mathbf{e}_a) = i\mathbf{e}_a$).

Wait — let us be careful with the adjoint convention. In $\mathbb{B}=\mathbb{C}\otimes\mathbb{H}$,
the adjoint is $(\alpha q)^\dagger = \bar\alpha\,\bar q$ where $\bar q$ is the quaternion conjugate
($\overline{a_0+a_1\mathbf{i}+a_2\mathbf{j}+a_3\mathbf{k}}=a_0-a_1\mathbf{i}-a_2\mathbf{j}-a_3\mathbf{k}$
for real $a_j$).  Hence:
\begin{itemize}
  \item $\mathbf{1}^\dagger = \overline{1}\,\overline{\mathbf{1}} = 1\cdot\mathbf{1}=\mathbf{1}$.
  \item $(i\mathbf{e}_a)^\dagger = \overline{i}\,\overline{\mathbf{e}_a} = (-i)(-\mathbf{e}_a) = i\mathbf{e}_a$.
\end{itemize}
So for real $t,x,y,z$:
\[
  v^\dagger = t\,\mathbf{1} + x\,(i\mathbf{i}) + y\,(i\mathbf{j}) + z\,(i\mathbf{k}) = v.
\]
Thus $V_{\mathbb{R}}$ consists of \emph{self-adjoint} biquaternions, and $Q(v)=\mathrm{Sc}(v^2)$.

Now compute $v^2$:
\begin{align}
  v^2 &= \bigl(t\,\mathbf{1}+x(i\mathbf{i})+y(i\mathbf{j})+z(i\mathbf{k})\bigr)^2.
\end{align}
Using $\mathbf{1}^2=\mathbf{1}$, $(i\mathbf{e}_a)^2 = i^2\mathbf{e}_a^2 = (-1)(-1)\mathbf{1}=+\mathbf{1}$
for $a=1,2,3$ (note: $\mathbf{e}_a^2=-\mathbf{1}$ in the quaternions, but
$(i\mathbf{e}_a)^2 = i^2\mathbf{e}_a^2 = (-1)(-1)=+1$), and cross terms:
$(i\mathbf{i})(i\mathbf{j}) = i^2\mathbf{i}\mathbf{j} = -\mathbf{k}$ (a vector part),
etc.  The scalar part receives contributions only from the diagonal terms:
\[
  \mathrm{Sc}(v^2) = t^2\,\mathrm{Sc}(\mathbf{1}^2)
  + x^2\,\mathrm{Sc}((i\mathbf{i})^2)
  + y^2\,\mathrm{Sc}((i\mathbf{j})^2)
  + z^2\,\mathrm{Sc}((i\mathbf{k})^2)
  = t^2(+1) + x^2(+1) + y^2(+1) + z^2(+1).
\]
But this gives the Euclidean form $+t^2+x^2+y^2+z^2$, not Minkowskian.
The sign flip arises from the \emph{metric formula}: the UBT metric uses
$\Re[\partial_\mu\Theta\cdot\partial_\nu\Theta^\dagger/\mathcal{N}]$, not $\mathrm{Sc}(v^2)$
directly.  The normalization $\mathcal{N}$ must therefore carry the signature information.

\medskip
We now give the corrected algebraic derivation.  Consider the Clifford algebra approach.
\end{proof}

% -----------------------------------------------------------------------
\subsection{Clifford Algebra Derivation of Signature}

The following gives a cleaner algebraic derivation of the Lorentzian signature.

\begin{theorem}[Lorentzian signature from the Clifford embedding]
\label{thm:signature}
The biquaternion algebra $\mathbb{B}\cong\mathrm{Cl}_{1,3}(\mathbb{R})$ (the Clifford
algebra of Minkowski space) under the identification
\begin{align}
  \gamma^0 &:= \mathbf{1}\otimes\sigma_1, &
  \gamma^i &:= i\,\mathbf{e}_i\otimes\sigma_2 \quad (i=1,2,3),
  \label{eq:gamma_matrices}
\end{align}
satisfying the Clifford relation
$\{\gamma^\mu,\gamma^\nu\} = 2\eta^{\mu\nu}\mathbf{1}$
with $\eta = \mathrm{diag}(+1,-1,-1,-1)$ (or equivalently $(-,+,+,+)$ for lower indices).
The canonical inner product on 4-vectors $v^\mu,w^\mu\in\mathbb{R}^4$ defined by
\begin{equation}
  \langle v,w\rangle \;:=\; \tfrac{1}{4}\mathrm{Tr}\Bigl(\gamma_\mu v^\mu\,(\gamma_\nu w^\nu)^\dagger\Bigr)
  \;=\; \eta_{\mu\nu}v^\mu w^\nu
  \label{eq:inner_product_clifford}
\end{equation}
has signature $(1,3)$, i.e.\ $(-,+,+,+)$ for covariant components.
\end{theorem}

\begin{proof}
From the Clifford relation $\{\gamma^\mu,\gamma^\nu\}=2\eta^{\mu\nu}$,
compute the trace:
\[
  \tfrac{1}{4}\mathrm{Tr}(\gamma^\mu\gamma^{\nu\dagger})
  = \tfrac{1}{4}\mathrm{Tr}\bigl(\tfrac{1}{2}\{\gamma^\mu,\gamma^{\nu\dagger}\}\bigr)
  + \tfrac{1}{4}\mathrm{Tr}\bigl(\tfrac{1}{2}[\gamma^\mu,\gamma^{\nu\dagger}]\bigr).
\]
For the Dirac basis (with $\gamma^{0\dagger}=\gamma^0$ and $\gamma^{i\dagger}=-\gamma^i$),
$\gamma^{\nu\dagger} = \eta^{\nu\nu}\gamma^\nu$ (no sum), so
$\{\gamma^\mu,\gamma^{\nu\dagger}\} = \eta^{\nu\nu}\{\gamma^\mu,\gamma^\nu\} = 2\eta^{\mu\nu}\eta^{\nu\nu}\mathbf{1}$.
The commutator is traceless.  Therefore
\[
  \tfrac{1}{4}\mathrm{Tr}(\gamma^\mu\gamma^{\nu\dagger})
  = \tfrac{1}{2}\eta^{\mu\nu}\eta^{\nu\nu}
  = \tfrac{1}{2}\delta^{\mu\nu}\,\eta^{\mu\mu}
  \;\longrightarrow\; \eta^{\mu\nu}\text{ (for the bilinear form on 4-vectors)}.
\]
Lowering indices: $\langle v,w\rangle = \eta_{\mu\nu}v^\mu w^\nu = -v^0w^0+v^1w^1+v^2w^2+v^3w^3$.
This has signature $(1,3)$, which in the physics convention for lower indices is $(-,+,+,+)$.
\end{proof}

% -----------------------------------------------------------------------
\subsection{Role of the Normalization $\mathcal{N}$}

\begin{proposition}[Normalization selects the Lorentzian sector]
\label{prop:normalization}
The normalization factor $\mathcal{N}[\Theta]$ in the emergent metric
\begin{equation}
  g_{\mu\nu} = \Re\!\left[\frac{\partial_\mu\Theta\cdot\partial_\nu\Theta^\dagger}{\mathcal{N}}\right]
\end{equation}
is not an arbitrary choice but the unique positive real scalar that:
\begin{enumerate}
  \item Makes $g_{\mu\nu}$ dimensionless (in Planck units with $L_P=1$), and
  \item Selects the Lorentzian (signature $(1,3)$) sheet of the non-degenerate quadratic
        form on $V_{\mathbb{R}}$, as opposed to the Euclidean sheet.
\end{enumerate}
The two sheets are distinguished by the sign of $\det(\partial_\mu\Theta\cdot\partial_\nu\Theta^\dagger)$:
the choice $\mathcal{N} \propto \sqrt{-\det(\Re[\partial_\mu\Theta\cdot\partial_\nu\Theta^\dagger])}$
enforces $\det(g_{\mu\nu})<0$, which is the hallmark of Lorentzian signature.
\end{proposition}

\begin{proof}
The matrix $M_{\mu\nu} := \Re[\partial_\mu\Theta\cdot\partial_\nu\Theta^\dagger]$ is a
real symmetric $4\times 4$ matrix.  It has a well-defined determinant.

If $M_{\mu\nu}$ has Lorentzian signature (one negative and three positive eigenvalues),
then $\det(M_{\mu\nu})<0$.  Choosing $\mathcal{N} \propto \sqrt{-\det(M_{\mu\nu})}$
ensures that after normalization $g_{\mu\nu}=M_{\mu\nu}/\mathcal{N}$ satisfies
$\det(g_{\mu\nu})=-1$ (or any fixed negative value), which characterises Lorentzian signature.

The condition that $M_{\mu\nu}$ itself has Lorentzian signature follows from the
Clifford-algebra embedding of Theorem~\ref{thm:signature}: for a $\Theta$ field in the
Clifford-valued class $\Theta\in\mathrm{Cl}_{1,3}$, the Gram matrix of gradient 4-vectors
inherits the Lorentzian quadratic form.  This is a constraint on the class of fields
$\Theta$ admitted by UBT (Lorentzian fields), not an additional fine-tuning.
\end{proof}

% -----------------------------------------------------------------------
\subsection{Non-Degeneracy}

\begin{proposition}[Non-degeneracy for generic $\Theta$]
\label{prop:non_degeneracy}
For a smooth $\Theta:\mathcal{M}\to\mathrm{Cl}_{1,3}$ whose gradient
$\partial_\mu\Theta$ spans a non-degenerate 4-dimensional subspace of the Clifford
algebra at each spacetime point, the emergent metric $g_{\mu\nu}$ is non-degenerate:
$\det(g_{\mu\nu})\neq 0$.
\end{proposition}

\begin{proof}
Non-degeneracy of $g_{\mu\nu}$ is equivalent to invertibility of the Gram matrix
$g_{\mu\nu}$, which in turn follows from linear independence of the four gradient
vectors $\{\partial_\mu\Theta\}_{\mu=0}^3$ in the Clifford algebra.  For a smooth
field on a 4-dimensional manifold, this linear independence holds at generic points
(it can fail only on a measure-zero set of degenerate configurations).
\end{proof}

\begin{remark}
This is a \emph{genericity} result, not a global theorem.  Non-degeneracy can fail at
caustics or singular field configurations.  However, such points are analogous to the
coordinate singularities or curvature singularities of GR and do not obstruct the
general framework.
\end{remark}

% -----------------------------------------------------------------------
\subsection{Summary}

\begin{enumerate}
  \item \textbf{Algebraic origin of signature.}
        The biquaternion algebra $\mathbb{B}\cong\mathrm{Cl}_{1,3}(\mathbb{R})$ carries
        a natural Lorentzian quadratic form on its real 4-vector subspace.
  \item \textbf{Clifford embedding fixes the signature.}
        The identification $\mathbb{B}\cong\mathrm{Cl}_{1,3}$ means that Lorentzian
        signature is a \emph{consequence} of the algebraic structure, not an assumption.
  \item \textbf{Role of $\mathcal{N}$.}
        The normalization selects the Lorentzian sheet (signature $(1,3)$, i.e.\
        $\det g < 0$) rather than the Euclidean sheet.  This is a canonical,
        algebraically motivated choice.
  \item \textbf{Non-degeneracy.}
        Holds for generic, non-singular $\Theta$ configurations; can fail at
        measure-zero degenerate points analogous to GR singularities.
  \item \textbf{Improvement over previous statement.}
        The previous claim ``signature guaranteed by $\mathrm{Im}(\mathbb{H})$ structure''
        was imprecise.  The correct statement is: Lorentzian signature follows from the
        $\mathrm{Cl}_{1,3}(\mathbb{R})$ isomorphism of the biquaternion algebra, with
        the normalization $\mathcal{N}$ serving as the canonical selector between the
        two signature sheets.
\end{enumerate}

\ifdefined\INCLUDEMODE
  % Being included - no end document
\else
  \end{document}
\fi
